\documentclass[10pt,a4paper]{amsart}
\usepackage[margin=19mm]{geometry}
\usepackage{amsmath,amssymb,mathtools}
\usepackage[scr=rsfs,cal=euler]{mathalfa}
\usepackage{xfrac}
\usepackage[unicode,hidelinks,bookmarks=false]{hyperref}
\newcommand{\ie}{i.e.\ }
\newcommand{\cf}{cf.\ }
\newcommand{\resp}{respectively\ }
\newcommand{\Subsection}{\subsection}
\providecommand{\nobreakdash}{\nobreak\hbox{-}}
\setlength{\emergencystretch}{2em}
\multlinegap0pt
\usepackage{amssymb}
\usepackage{tikz}
\usepackage{enumerate}
\newtheorem{prop}{Proposition}
\newcommand{\Hom}{\mathrm{Hom}}
\newcommand{\KK}{\mathbb{K}}
\newcommand{\Vect}{\mathrm{Vect}}
\newcommand{\ch}[1]{\mathrm{Ch^*({#1})}}
\title{Formality of hypercommutative algebras of Kähler\\ and Calabi--Yau manifolds}
\author{Joana Cirici, Geoffroy Horel}
\date{Source: arXiv:2302.08492v3 (CC BY 4.0)}
\begin{document}
\maketitle

\noindent\emph{Source and licence.} Formality of hypercommutative algebras of Kähler\\ and Calabi--Yau manifolds. Version arXiv:2302.08492v3 (CC BY 4.0), distributed by arXiv under the Creative Commons Attribution 4.0 International licence, http://creativecommons.org/licenses/by/4.0/, as recorded in the arXiv licence metadata for this exact version and shown on the abstract page https://arxiv.org/abs/2302.08492v3. Full licence notice: \texttt{licenses/arxiv-2302.08492-CC-BY-4.0.txt}.\par\smallskip

\noindent\textit{Editorial scope.} Counted unit: the article's prop prop (source0.tex lines 662--666). The article's complete original proof of it is delivered with it (source0.tex lines 668--674, one \texttt{proof} environment of 7 lines). No mathematics is written, repaired or re-derived: the statement and the proof are the article's own lines, in the article's own notation, and no retained line is substituted or re-flowed.  No further source-local context is needed: every cross-reference of every retained line resolves inside the two delivered spans.  Nothing was omitted from any retained span, so the package carries no omission record.  No retained line contains a citation, so the wrapper carries no bibliography and no bibliography entry.  No retained line contains a figure environment, an image-inclusion command or drawing code, so no image is delivered and the package declares no asset dependency.  Editorial formatting only, in addition: an AMS \texttt{amsart} preamble that restates the article's own declarations and macros used by the retained lines, namely the environments \texttt{prop} on separate counters, together with the standard packages those lines need.  The retained environments print in this document as Proposition 1 in order of appearance, whereas the article prints the same items with the numbering of its own full document.  The complete native source of the article as distributed in the arXiv e-print for this exact version is bundled unchanged as \texttt{sources/137488/source0.tex}.
\begin{prop}
Let $\alpha=r/s$ with $r$ and $s$ two coprime integers. The cohomology functor
\[\ch{Gr\Vect_\KK}^{\alpha\text{-pure}}\to gr\Vect_\KK\]
induces a symmetric monoidal equivalence of $\infty$-categories onto the full subcategory of $gr\Vect_\KK$ spanned by graded vector spaces that are zero except in degree divisible by $s$.
\end{prop}

\begin{proof}
This functor is clearly essentially surjective. In order to prove that it is fully faithful, it suffices to observe the following formula for the homotopy groups of the mapping spaces in $\ch{Gr\Vect_\KK}$
\[\pi_k\mathrm{Map}_{\ch{Gr\Vect_\KK}}(C,D)=\bigoplus_{i\in\mathbb{Z}}[C_i,D_i[k]]\]
in particular, we see that if $C$ and $D$ are both $\alpha$-pure, all the higher homotopy groups are trivial and we have 
\[\pi_0\mathrm{Map}_{\ch{Gr\Vect_\KK}}(C,D)=\bigoplus_{j\in\mathbb{Z}}\Hom(H^{sj}(C_{rj}),H^{sj}(D_{rj}))\]
This shows that the cohomology functor is fully faithful.
\end{proof}

\end{document}
